\section{Conclusion}
\label{sec:conclusion}

We presented a consistent continuum formulation of a deterministic substrate theory in which
observable physics arises from waves on a rotationally invariant 3D elastic medium embedded in $\R^4$,
with external time as evolution parameter.
The continuum model was stated as an isotropic hyperelastic action built from the induced metric of the embedding,
including a simple reference-metric mismatch parameter $\alpha$ that encodes homogeneous pre-stress.

Assuming an ordered narrowband carrier sector, we constructed the Berry (rank-1) and Wilczek--Zee (non-Abelian)
connections induced by polarization transport and interpreted them as the natural origin of gauge degrees of freedom
in a coarse-grained description.
We then outlined how localized, particle-like states can be posed as self-adjoint eigenproblems around symmetric
backgrounds, with angular structure organized by spherical-harmonic mode families and confinement attributed to
geometric/constitutive nonlinearity (self-guidance) rather than to fundamental point charges.

The resulting picture is a theory fragment that ties together:
(1) a continuum brane substrate,
(2) an ordered narrowband phase sector,
(3) geometric connection/holonomy as gauge structure,
and (4) solitonic fermions with spinorial ($4\pi$) holonomy attributed to Berry/Wilczek--Zee transport.
Discrete numerical realizations are treated only as implementation reference material and are placed in an appendix.
